\section{Nontrivial character labels with one fixed zeta denominator}
\label{sec:sp4-fixed-spectrum}

The fourfold symplectic detector has a different family parameter from
the cubic character extraction in the preceding sections. One can vary
the inducing character without varying the target zeta function. The
companion conductors consequently remain fixed. We give the exact
application and its limitations. Its algebraic character statements and
finite-sum obstruction are kernel checked; its automorphic realization
uses \cite[Section 6.5]{GaoInverseEis} and the explicitly stated local
analytic interface in Section~\ref{sec:sp4-principal-window}.

\paragraph{The complete kernel of the two root restrictions.}
Let \(A,C\) be commutative groups. A character \(f:A\times A\to C\)
satisfies
\[
 f(a,a^{-1})=1,\qquad f(1,a^2)=1\quad(a\in A)
\]
if and only if
\[
 f(a,b)=\kappa(ab),\qquad \kappa(a)^2=1\quad(a\in A)
\]
for a character \(\kappa:A\to C\). Indeed set
\(\kappa(a)=f(1,a)\). The second restriction makes it quadratic.
The factorization
\((a,b)=(a,a^{-1})(1,ab)\) proves the formula. Conversely that
formula gives both restrictions directly. This is
\texttt{Metaplectic.Sp4FixedSpectrum.fixed\_roots\_iff\_quadratic}.

For two characters \(\tau,\sigma\), their determinant twists agree
on every square torus coordinate precisely when
\[
 \tau(a^2b^2)=\sigma(a^2b^2)\ (a,b\in A)
 \quad\Longleftrightarrow\quad \tau^2=\sigma^2.
\]
Set \(b=1\) for necessity and multiply the two equalities for
sufficiency. This equivalence is maintained as
\texttt{square\_torus\_twist\_eq\_iff} in the same namespace.
The theorem \texttt{quartic\_twist\_fixed\_roots} specializes the
first characterization to \(\kappa=\tau^2\) when \(\tau^4=1\).
These are abstract group statements; they do not assert that every
quadratic Hecke character has a quartic root.

\paragraph{A concrete infinite automorphic family.}
Fix \(K=\mathbb Q(\mu_8)\) and the principal genuine datum
\(\chi_0\) of Section~\ref{sec:sp4-principal-window}.
Twist its Siegel Levi \(M=\operatorname{GL}_2\) by
\(\tau\circ\det_M\), where \(\tau\) is a finite-order Hecke
character of \(K\) with \(\tau^4=1\). On the two lifted roots the
new factors are
\[
 \tau(a^2a^{-2})=1,\qquad \tau(a^4)=1.
\]
Thus the exceptional Levi value and the target character \(\eta=1\)
are both unchanged. On the central square coordinates the new factor
is \(\tau(a^2b^2)=\tau^2(ab)\). The labels are therefore the
quadratic characters \(\kappa=\tau^2\) that actually occur, rather
than the two choices of a quartic root. This is not a determinant
twist of the ambient symplectic group, whose determinant is one.

For an explicit unbounded family take rational primes \(q\equiv1\pmod8\),
\(q\ge17\), primitive quartic Dirichlet characters \(\rho_q\), and
\(\tau_q=\rho_q\circ N_{K/\mathbb Q}\). The square is nontrivial:
the quadratic extension of conductor \(q\) is disjoint from the field
of conductor \(8\), so its character does not become trivial on
restriction to \(K\). Different primes give different ramification
and hence different labels. Conjugate quartic roots at one prime
give the same square-coordinate label. No neutral label is required
to detect a zeta zero in this family.

Let \(S_0\) contain the archimedean and dyadic places and fix the
source measures and additive character. With \(u=4s+2\), Gao's
factorization for these data gives
\[
 E_{\tau,\psi}(s)=\frac{B_\tau(s)}{\zeta_K^{S_q}(u)},
 \qquad S_q=S_0\cup\{v:v\mid q\}.
\]
At a ramified finite place choose a Levi vector with Whittaker
functional one, and an open-cell section supported on
\(P w_\ell^{-1}N'(\mathcal O_v)\). On \(w_\ell^{-1}t\) its value
is that vector times \(\mathbf1_{N'(\mathcal O_v)}(t)\).
For unramified \(\psi_v\) and volume one the local Jacquet integral
is exactly one, independently of \(s\). A sufficiently small compact
subgroup with the appropriate volume normalization handles a fixed
bad place. Such a compact-modulo-\(P\) section is smooth and finite
under a suitable compact open subgroup. This is an existence
statement; its level and Sobolev norms are not bounded here.

The finite-order characters at the complex places are trivial, so
one fixed admissible archimedean section can be used for all \(q\).
Under the meromorphic continuity and dense finite-observation
hypotheses of Section~\ref{sec:sp4-principal-window}, choose it so
that its factor \(B(s)\) is holomorphic and nonzero at the specified
candidate zero. Normalize the finite local factors as above. Then
\(B_\tau=B\), independent of \(q\). This is finite-point
nonvanishing, not a uniform height or level estimate.

\paragraph{Restore the ramified factors before differentiating.}
Each such \(q\) splits into four degree-one places of \(K\). Hence
\begin{align*}
 P_q(s)&=(1-q^{-u})^4,&
 Z_\tau(s)&=P_q(s)E_{\tau,\psi}(s)
             =\frac{B(s)}{\zeta_K^{S_0}(u)},\\
 P_q'(s)&=16(\log q)q^{-u}(1-q^{-u})^3,&
 Z_\tau'(s)&=P_qE_{\tau,\psi}'+P_q'E_{\tau,\psi}.
\end{align*}
Both differentiated terms are required. Put
\[
 C(u)=L(u,\chi_{-4})L(u,\chi_8)L(u,\chi_{-8}).
\]
The factorization over the fixed field gives the exact meromorphic
identity
\begin{equation}\label{eq:sp4-fixed-spectrum-inverse}
 C(u)Z_\tau(s)=\frac{B(s)}{\zeta_{\mathbb Q}^{(2)}(u)}.
\end{equation}
The notation \((2)\) removes the Euler factor at two. A shared zero
of the companions causes no loss: its multiplicity cancels on both
sides, leaving the multiplicity of the rational zeta zero. Since
\(1-2^{-u}\ne0\) for \(\Re u>0\), this finite removal does not
change any nontrivial zero in that region. In particular a zero
\(\rho\) gives a pole at \(s_\rho=(\rho-2)/4\) whenever
\(B(s_\rho)\ne0\), for every nontrivial label above.

The full signed logarithmic identity is
\begin{multline}\label{eq:sp4-fixed-spectrum-log}
 \zeta_K^{S_0}(u)
   \bigl(P_qE_{\tau,\psi}'+P_q'E_{\tau,\psi}\bigr)-B'(s)
    +4B(s)\frac{C'(u)}{C(u)}\\
 =-4B(s)\frac{(\zeta_{\mathbb Q}^{(2)})'(u)}
                         {\zeta_{\mathbb Q}^{(2)}(u)}.
\end{multline}
This follows by differentiating \(Z_\tau=B/\zeta_K^{S_0}\)
and using \((\zeta_K^{S_0})'/\zeta_K^{S_0}
=(\zeta_{\mathbb Q}^{(2)})'/\zeta_{\mathbb Q}^{(2)}+C'/C\).
Every reciprocal in \(C'/C\), the derivative of \(B\), the chain
factor four and the principal signal remain in the identity.

The inverse consumer is simpler than a moving-character product:
its three companions have fixed conductors \(4,8,8\).
The already proved periodic-companion continuation estimate gives
\[
 |L(u,\chi)|\le |u|
   \left(\frac1{|u-1|}+\frac{2m+1}{\Re u}\right)
 \quad(\Re u>0,\ u\ne1)
\]
for a character modulo \(m\). On
\(1/2\le\Re u\le7/8\), it follows that
\(|C(u)|\le26\cdot42^2|u|^3=45864|u|^3\).
This specialization is a mathematical consequence of the compiled
provider, not an additional compiled declaration. There is no
growing conductor in this companion bound. The moving inducing
sections and their norms still have to be estimated independently.

\paragraph{Complete physical support and exponent conversion.}
Let \(a_K(n)\) be the odd-ideal coefficients of
\(1/\zeta_K^{S_0}(u)\). For every compactly supported smooth \(w\),
the exact coefficient cancellation is
\[
 \sum_{\substack{a,b_1,b_2,b_3\ge1\\ab_1b_2b_3\ {\rm odd}}}
 a_K(a)\chi_{-4}(b_1)\chi_8(b_2)\chi_{-8}(b_3)
       w(ab_1b_2b_3/X)
 =\sum_{\substack{n\ge1\\n\ {\rm odd}}}\mu(n)w(n/X).
\]
Both sides are finite on the displayed common product support.
There is no additional coprimality restriction, and all unit rows
are included. The ramified restoration has all five terms
\[
 P_q(s)=\sum_{h=0}^{4}(-1)^h\binom4h q^{-2h}(q^{4h})^{-s}.
\]
In a Mellin window of scale \(Y\) these give scales
\(Y/q^{4h}\); differentiation additionally gives the factors
\(-4h\log q\). A proposed bound on the whole packet must cover
these shifted scales, including small arguments.

For a finite-height integral let
\[
 Q_{\tau,w}(Y;T)=\frac1{2\pi i}
  \int_{c'-iT/4}^{c'+iT/4}
    C(4s+2)P_q(s)E_{\tau,\psi}(s)
    \widehat w(4s+2)Y^s\,ds,
 \qquad c'=(c-2)/4,\ c>1.
\]
The exact substitution \(u=4s+2\) gives
\[
 \frac1{2\pi i}\int_{c-iT}^{c+iT}
    \frac{B((u-2)/4)\widehat w(u)}{\zeta_{\mathbb Q}^{(2)}(u)}
       X^u\,du
       =4X^2 Q_{\tau,w}(X^4;T).
\]
Passage to an infinite contour requires the usual decay and horizontal
controls. If an independently proved full geometric period estimate
and those controls give \(Q(Y)\ll_\epsilon Y^{r+\epsilon}\),
with a detecting window at each candidate zero, the resulting
zero-free boundary is \(2+4r\). Thus the installed boundary \(7/8\)
corresponds to \(r=-9/32\); the target \(3/4\) requires
\(r=-5/16\), and \(1/2\) requires \(r=-3/8\).
No such independent improved period estimate is supplied here.
For \(q\asymp X^\delta\), an additional cost \(q^A\), length
replacement \(Y\mapsto Yq^\lambda\), and normalization \(q^{-\nu}\)
would change the boundary to
\(2+4r+\delta(A+\lambda r-\nu)\); none may be omitted.

\paragraph{A proved obstruction to the wrong physical kernel.}
For finite fields \(k\), any multiplicative character \(\chi\),
and any weight \(w:k\to\mathbb C\), set
\[
 R_\chi(w)=\sum_{a,b,c\in k}\chi(ac-b^2)w(c).
\]
If a unit \(v\) has \(\chi(v)^2\ne1\), the bijection
\((a,b)\mapsto(v^2a,vb)\) gives
\(R_\chi(w)=\chi(v)^2R_\chi(w)\), hence \(R_\chi(w)=0\).
The proof keeps the entire finite sum and is maintained as
\texttt{Metaplectic.Sp4FixedSpectrum.rank\_one\_probe\_eq\_zero}.
Consequently \(\chi(\det B)w(B_{22})\) alone cannot represent the
nonzero rank-one automorphic observation. The full Levi Whittaker
factor cannot be removed. The existing quadratic reciprocity-energy
and sextic Eisenstein providers do not identify that factor: the
latter also has different base field and row hypotheses.

Nontrivial labels retain a common zero signal, but square-frequency
terms in their arithmetic kernels can still survive. Moreover the
normalized periods in \eqref{eq:sp4-fixed-spectrum-inverse} are
identical: any finite linear combination with total coefficient zero
erases the signal. A useful joint estimate must keep total coefficient
one and bound the full geometric period with its moving section norms.
This is a distinct consumer with fixed companions, rather than a
transfer of the cubic \(X^{3/16}\) error. The certified global
boundary remains \(7/8\). Priority for this combined application is
unestablished; the supporting character arguments are not claimed novel.
